\documentclass[11pt]{article}
% jugeo-common.tex — shared preamble for all Judgment Geometry papers
% Include via: % jugeo-common.tex — shared preamble for all Judgment Geometry papers
% Include via: % jugeo-common.tex — shared preamble for all Judgment Geometry papers
% Include via: \input{jugeo-common}

% ── Packages ────────────────────────────────────────────────────────────
\usepackage{amsmath,amssymb,amsthm,mathtools}
\usepackage{tikz-cd}
\usepackage{listings}
\usepackage{xcolor}
\usepackage{xspace}
\usepackage{booktabs}
\usepackage{multirow}
\usepackage{enumitem}
\usepackage[expansion=false]{microtype}
\usepackage{hyperref}
\usepackage{cleveref}
% \usepackage{thmtools}  % disabled: not available in this TeX installation
\usepackage{float}
\usepackage{caption}
\usepackage{subcaption}
\usepackage{tabularx}
\usepackage{stmaryrd}       % \llbracket, \rrbracket
\usepackage{bbm}            % \mathbbm{1}

% ── Page geometry ───────────────────────────────────────────────────────
\usepackage[margin=1in]{geometry}

% ── Theorem environments ────────────────────────────────────────────────
\theoremstyle{plain}
\newtheorem{theorem}{Theorem}[section]
\newtheorem{lemma}[theorem]{Lemma}
\newtheorem{proposition}[theorem]{Proposition}
\newtheorem{corollary}[theorem]{Corollary}

\theoremstyle{definition}
\newtheorem{definition}[theorem]{Definition}
\newtheorem{example}[theorem]{Example}
\newtheorem{construction}[theorem]{Construction}

\theoremstyle{remark}
\newtheorem{remark}[theorem]{Remark}
\newtheorem{notation}[theorem]{Notation}

% ── Python listings style ──────────────────────────────────────────────
\definecolor{jugeoBlue}{HTML}{2563EB}
\definecolor{jugeoGreen}{HTML}{059669}
\definecolor{jugeoRed}{HTML}{DC2626}
\definecolor{jugeoPurple}{HTML}{7C3AED}
\definecolor{jugeoGray}{HTML}{6B7280}
\definecolor{jugeoBg}{HTML}{F8FAFC}

\lstdefinestyle{jugeo-python}{
  language=Python,
  basicstyle=\ttfamily\small,
  keywordstyle=\color{jugeoBlue}\bfseries,
  stringstyle=\color{jugeoGreen},
  commentstyle=\color{jugeoGray}\itshape,
  numberstyle=\tiny\color{jugeoGray},
  backgroundcolor=\color{jugeoBg},
  numbers=left,
  numbersep=6pt,
  frame=single,
  framerule=0.4pt,
  rulecolor=\color{jugeoGray!40},
  breaklines=true,
  breakatwhitespace=true,
  showstringspaces=false,
  tabsize=4,
  xleftmargin=1.5em,
  framexleftmargin=1.5em,
  morekeywords={dataclass,frozen,field,Enum,IntEnum,auto,Any,
                Mapping,Sequence,Optional,match,case},
  literate={→}{{$\to$}}1 {←}{{$\leftarrow$}}1
           {≤}{{$\leq$}}1 {≥}{{$\geq$}}1 {≠}{{$\neq$}}1
           {∈}{{$\in$}}1 {∀}{{$\forall$}}1 {∃}{{$\exists$}}1
           {⊕}{{$\oplus$}}1 {⊖}{{$\ominus$}}1,
  escapeinside={(*@}{@*)},
}
\lstset{style=jugeo-python}

\lstdefinestyle{jugeo-output}{
  basicstyle=\ttfamily\small,
  backgroundcolor=\color{jugeoBg},
  frame=single,
  framerule=0.4pt,
  rulecolor=\color{jugeoGray!40},
  breaklines=true,
  showstringspaces=false,
  xleftmargin=1.5em,
  framexleftmargin=0.5em,
  numbers=none,
}

% ── JuGeo-specific macros ──────────────────────────────────────────────

% Judgment 8-tuple
\newcommand{\J}{\mathcal{J}}
\newcommand{\Jud}[8]{(#1,\,#2,\,#3,\,#4,\,#5,\,#6,\,#7,\,#8)}
\newcommand{\JudShort}[1]{\J_{#1}}

% Components
\newcommand{\coord}{\mathsf{c}}            % coordinate
\newcommand{\prop}{\varphi}                 % proposition
\newcommand{\carrier}{\mathcal{A}}          % carrier type
\newcommand{\evid}{\mathcal{E}}             % evidence bundle
\newcommand{\oblig}{\mathcal{O}}            % obligations
\newcommand{\obstr}{\mathcal{B}}            % obstructions
\newcommand{\trust}{\mathcal{T}}            % trust annotation
\newcommand{\prov}{\Pi}                     % provenance

% Trust algebra
\newcommand{\Talg}{\mathfrak{T}}            % trust ordered algebra
\newcommand{\Eadm}{\mathcal{E}_{\mathrm{adm}}}
\newcommand{\tleq}{\preceq}                 % trust ordering
\newcommand{\tjoin}{\oplus}                 % conservative join
\newcommand{\tatten}{\ominus}               % attenuation
\newcommand{\tpromote}{\uparrow_\pi}        % promotion
\newcommand{\tdemote}{\downarrow_\chi}       % demotion

% Trust tiers
\newcommand{\Tcontra}{\ensuremath{\mathsf{CONTRADICTED}}}
\newcommand{\Tunver}{\ensuremath{\mathsf{UNVERIFIED}}}
\newcommand{\Tcopilot}{\ensuremath{\mathsf{COPILOT}}}
\newcommand{\Toracle}{\ensuremath{\mathsf{ORACLE}}}
\newcommand{\Truntime}{\ensuremath{\mathsf{RUNTIME}}}
\newcommand{\Tsolver}{\ensuremath{\mathsf{SOLVER}}}
\newcommand{\Tproof}{\ensuremath{\mathsf{PROOF}}}

% Site and geometry
\newcommand{\Site}{\mathcal{S}}             % semantic site
\newcommand{\Cov}{\mathrm{Cov}}            % covering family
\newcommand{\Ob}{\mathrm{Ob}}              % objects
\newcommand{\Mor}{\mathrm{Mor}}            % morphisms
\newcommand{\res}{\mathrm{res}}            % restriction
\newcommand{\Cech}{\check{C}}              % Čech complex
\newcommand{\Hcoh}[1]{H^{#1}}             % cohomology group

% Descent
\newcommand{\descent}{\mathrm{desc}}
\newcommand{\glue}{\mathrm{glue}}
\newcommand{\overlap}[2]{#1 \cap #2}

% Semantic moves
\newcommand{\VERIFY}{\textsc{Verify}}
\newcommand{\CONSTRUCT}{\textsc{Construct}}
\newcommand{\NORMALIZE}{\textsc{Normalize}}
\newcommand{\GLUE}{\textsc{Glue}}
\newcommand{\RESTRICT}{\textsc{Restrict}}
\newcommand{\TRANSPORT}{\textsc{Transport}}
\newcommand{\REPAIR}{\textsc{Repair}}
\newcommand{\PROMOTE}{\textsc{Promote}}
\newcommand{\CHALLENGE}{\textsc{Challenge}}
\newcommand{\ESCALATE}{\textsc{Escalate}}

% Fragments
\newcommand{\QFLIA}{\texttt{QF\_LIA}}
\newcommand{\QFLRA}{\texttt{QF\_LRA}}
\newcommand{\QFBV}{\texttt{QF\_BV}}
\newcommand{\QFUF}{\texttt{QF\_UF}}

% Misc
\newcommand{\Python}{\textsc{Python}}
\newcommand{\JuGeo}{\textsc{JuGeo}}
\newcommand{\LEAN}{\textsc{Lean}\,4}
\newcommand{\Fstar}{F$^{\star}$}
\newcommand{\Coq}{\textsc{Coq}}
\newcommand{\Dafny}{\textsc{Dafny}}

% Common shortcuts
\newcommand{\ie}{i.e.\@\xspace}
\newcommand{\eg}{e.g.\@\xspace}
\newcommand{\cf}{cf.\@\xspace}
\newcommand{\wrt}{w.r.t.\@\xspace}

% Evaluation shortcuts
\newcommand{\numcases}{300}
\newcommand{\numspec}{100}
\newcommand{\numeq}{100}
\newcommand{\numbug}{100}
\newcommand{\medtime}{6.5\,ms}
\newcommand{\totaltime}{1.949\,s}


% ── Packages ────────────────────────────────────────────────────────────
\usepackage{amsmath,amssymb,amsthm,mathtools}
\usepackage{tikz-cd}
\usepackage{listings}
\usepackage{xcolor}
\usepackage{xspace}
\usepackage{booktabs}
\usepackage{multirow}
\usepackage{enumitem}
\usepackage[expansion=false]{microtype}
\usepackage{hyperref}
\usepackage{cleveref}
% \usepackage{thmtools}  % disabled: not available in this TeX installation
\usepackage{float}
\usepackage{caption}
\usepackage{subcaption}
\usepackage{tabularx}
\usepackage{stmaryrd}       % \llbracket, \rrbracket
\usepackage{bbm}            % \mathbbm{1}

% ── Page geometry ───────────────────────────────────────────────────────
\usepackage[margin=1in]{geometry}

% ── Theorem environments ────────────────────────────────────────────────
\theoremstyle{plain}
\newtheorem{theorem}{Theorem}[section]
\newtheorem{lemma}[theorem]{Lemma}
\newtheorem{proposition}[theorem]{Proposition}
\newtheorem{corollary}[theorem]{Corollary}

\theoremstyle{definition}
\newtheorem{definition}[theorem]{Definition}
\newtheorem{example}[theorem]{Example}
\newtheorem{construction}[theorem]{Construction}

\theoremstyle{remark}
\newtheorem{remark}[theorem]{Remark}
\newtheorem{notation}[theorem]{Notation}

% ── Python listings style ──────────────────────────────────────────────
\definecolor{jugeoBlue}{HTML}{2563EB}
\definecolor{jugeoGreen}{HTML}{059669}
\definecolor{jugeoRed}{HTML}{DC2626}
\definecolor{jugeoPurple}{HTML}{7C3AED}
\definecolor{jugeoGray}{HTML}{6B7280}
\definecolor{jugeoBg}{HTML}{F8FAFC}

\lstdefinestyle{jugeo-python}{
  language=Python,
  basicstyle=\ttfamily\small,
  keywordstyle=\color{jugeoBlue}\bfseries,
  stringstyle=\color{jugeoGreen},
  commentstyle=\color{jugeoGray}\itshape,
  numberstyle=\tiny\color{jugeoGray},
  backgroundcolor=\color{jugeoBg},
  numbers=left,
  numbersep=6pt,
  frame=single,
  framerule=0.4pt,
  rulecolor=\color{jugeoGray!40},
  breaklines=true,
  breakatwhitespace=true,
  showstringspaces=false,
  tabsize=4,
  xleftmargin=1.5em,
  framexleftmargin=1.5em,
  morekeywords={dataclass,frozen,field,Enum,IntEnum,auto,Any,
                Mapping,Sequence,Optional,match,case},
  literate={→}{{$\to$}}1 {←}{{$\leftarrow$}}1
           {≤}{{$\leq$}}1 {≥}{{$\geq$}}1 {≠}{{$\neq$}}1
           {∈}{{$\in$}}1 {∀}{{$\forall$}}1 {∃}{{$\exists$}}1
           {⊕}{{$\oplus$}}1 {⊖}{{$\ominus$}}1,
  escapeinside={(*@}{@*)},
}
\lstset{style=jugeo-python}

\lstdefinestyle{jugeo-output}{
  basicstyle=\ttfamily\small,
  backgroundcolor=\color{jugeoBg},
  frame=single,
  framerule=0.4pt,
  rulecolor=\color{jugeoGray!40},
  breaklines=true,
  showstringspaces=false,
  xleftmargin=1.5em,
  framexleftmargin=0.5em,
  numbers=none,
}

% ── JuGeo-specific macros ──────────────────────────────────────────────

% Judgment 8-tuple
\newcommand{\J}{\mathcal{J}}
\newcommand{\Jud}[8]{(#1,\,#2,\,#3,\,#4,\,#5,\,#6,\,#7,\,#8)}
\newcommand{\JudShort}[1]{\J_{#1}}

% Components
\newcommand{\coord}{\mathsf{c}}            % coordinate
\newcommand{\prop}{\varphi}                 % proposition
\newcommand{\carrier}{\mathcal{A}}          % carrier type
\newcommand{\evid}{\mathcal{E}}             % evidence bundle
\newcommand{\oblig}{\mathcal{O}}            % obligations
\newcommand{\obstr}{\mathcal{B}}            % obstructions
\newcommand{\trust}{\mathcal{T}}            % trust annotation
\newcommand{\prov}{\Pi}                     % provenance

% Trust algebra
\newcommand{\Talg}{\mathfrak{T}}            % trust ordered algebra
\newcommand{\Eadm}{\mathcal{E}_{\mathrm{adm}}}
\newcommand{\tleq}{\preceq}                 % trust ordering
\newcommand{\tjoin}{\oplus}                 % conservative join
\newcommand{\tatten}{\ominus}               % attenuation
\newcommand{\tpromote}{\uparrow_\pi}        % promotion
\newcommand{\tdemote}{\downarrow_\chi}       % demotion

% Trust tiers
\newcommand{\Tcontra}{\ensuremath{\mathsf{CONTRADICTED}}}
\newcommand{\Tunver}{\ensuremath{\mathsf{UNVERIFIED}}}
\newcommand{\Tcopilot}{\ensuremath{\mathsf{COPILOT}}}
\newcommand{\Toracle}{\ensuremath{\mathsf{ORACLE}}}
\newcommand{\Truntime}{\ensuremath{\mathsf{RUNTIME}}}
\newcommand{\Tsolver}{\ensuremath{\mathsf{SOLVER}}}
\newcommand{\Tproof}{\ensuremath{\mathsf{PROOF}}}

% Site and geometry
\newcommand{\Site}{\mathcal{S}}             % semantic site
\newcommand{\Cov}{\mathrm{Cov}}            % covering family
\newcommand{\Ob}{\mathrm{Ob}}              % objects
\newcommand{\Mor}{\mathrm{Mor}}            % morphisms
\newcommand{\res}{\mathrm{res}}            % restriction
\newcommand{\Cech}{\check{C}}              % Čech complex
\newcommand{\Hcoh}[1]{H^{#1}}             % cohomology group

% Descent
\newcommand{\descent}{\mathrm{desc}}
\newcommand{\glue}{\mathrm{glue}}
\newcommand{\overlap}[2]{#1 \cap #2}

% Semantic moves
\newcommand{\VERIFY}{\textsc{Verify}}
\newcommand{\CONSTRUCT}{\textsc{Construct}}
\newcommand{\NORMALIZE}{\textsc{Normalize}}
\newcommand{\GLUE}{\textsc{Glue}}
\newcommand{\RESTRICT}{\textsc{Restrict}}
\newcommand{\TRANSPORT}{\textsc{Transport}}
\newcommand{\REPAIR}{\textsc{Repair}}
\newcommand{\PROMOTE}{\textsc{Promote}}
\newcommand{\CHALLENGE}{\textsc{Challenge}}
\newcommand{\ESCALATE}{\textsc{Escalate}}

% Fragments
\newcommand{\QFLIA}{\texttt{QF\_LIA}}
\newcommand{\QFLRA}{\texttt{QF\_LRA}}
\newcommand{\QFBV}{\texttt{QF\_BV}}
\newcommand{\QFUF}{\texttt{QF\_UF}}

% Misc
\newcommand{\Python}{\textsc{Python}}
\newcommand{\JuGeo}{\textsc{JuGeo}}
\newcommand{\LEAN}{\textsc{Lean}\,4}
\newcommand{\Fstar}{F$^{\star}$}
\newcommand{\Coq}{\textsc{Coq}}
\newcommand{\Dafny}{\textsc{Dafny}}

% Common shortcuts
\newcommand{\ie}{i.e.\@\xspace}
\newcommand{\eg}{e.g.\@\xspace}
\newcommand{\cf}{cf.\@\xspace}
\newcommand{\wrt}{w.r.t.\@\xspace}

% Evaluation shortcuts
\newcommand{\numcases}{300}
\newcommand{\numspec}{100}
\newcommand{\numeq}{100}
\newcommand{\numbug}{100}
\newcommand{\medtime}{6.5\,ms}
\newcommand{\totaltime}{1.949\,s}


% ── Packages ────────────────────────────────────────────────────────────
\usepackage{amsmath,amssymb,amsthm,mathtools}
\usepackage{tikz-cd}
\usepackage{listings}
\usepackage{xcolor}
\usepackage{xspace}
\usepackage{booktabs}
\usepackage{multirow}
\usepackage{enumitem}
\usepackage[expansion=false]{microtype}
\usepackage{hyperref}
\usepackage{cleveref}
% \usepackage{thmtools}  % disabled: not available in this TeX installation
\usepackage{float}
\usepackage{caption}
\usepackage{subcaption}
\usepackage{tabularx}
\usepackage{stmaryrd}       % \llbracket, \rrbracket
\usepackage{bbm}            % \mathbbm{1}

% ── Page geometry ───────────────────────────────────────────────────────
\usepackage[margin=1in]{geometry}

% ── Theorem environments ────────────────────────────────────────────────
\theoremstyle{plain}
\newtheorem{theorem}{Theorem}[section]
\newtheorem{lemma}[theorem]{Lemma}
\newtheorem{proposition}[theorem]{Proposition}
\newtheorem{corollary}[theorem]{Corollary}

\theoremstyle{definition}
\newtheorem{definition}[theorem]{Definition}
\newtheorem{example}[theorem]{Example}
\newtheorem{construction}[theorem]{Construction}

\theoremstyle{remark}
\newtheorem{remark}[theorem]{Remark}
\newtheorem{notation}[theorem]{Notation}

% ── Python listings style ──────────────────────────────────────────────
\definecolor{jugeoBlue}{HTML}{2563EB}
\definecolor{jugeoGreen}{HTML}{059669}
\definecolor{jugeoRed}{HTML}{DC2626}
\definecolor{jugeoPurple}{HTML}{7C3AED}
\definecolor{jugeoGray}{HTML}{6B7280}
\definecolor{jugeoBg}{HTML}{F8FAFC}

\lstdefinestyle{jugeo-python}{
  language=Python,
  basicstyle=\ttfamily\small,
  keywordstyle=\color{jugeoBlue}\bfseries,
  stringstyle=\color{jugeoGreen},
  commentstyle=\color{jugeoGray}\itshape,
  numberstyle=\tiny\color{jugeoGray},
  backgroundcolor=\color{jugeoBg},
  numbers=left,
  numbersep=6pt,
  frame=single,
  framerule=0.4pt,
  rulecolor=\color{jugeoGray!40},
  breaklines=true,
  breakatwhitespace=true,
  showstringspaces=false,
  tabsize=4,
  xleftmargin=1.5em,
  framexleftmargin=1.5em,
  morekeywords={dataclass,frozen,field,Enum,IntEnum,auto,Any,
                Mapping,Sequence,Optional,match,case},
  literate={→}{{$\to$}}1 {←}{{$\leftarrow$}}1
           {≤}{{$\leq$}}1 {≥}{{$\geq$}}1 {≠}{{$\neq$}}1
           {∈}{{$\in$}}1 {∀}{{$\forall$}}1 {∃}{{$\exists$}}1
           {⊕}{{$\oplus$}}1 {⊖}{{$\ominus$}}1,
  escapeinside={(*@}{@*)},
}
\lstset{style=jugeo-python}

\lstdefinestyle{jugeo-output}{
  basicstyle=\ttfamily\small,
  backgroundcolor=\color{jugeoBg},
  frame=single,
  framerule=0.4pt,
  rulecolor=\color{jugeoGray!40},
  breaklines=true,
  showstringspaces=false,
  xleftmargin=1.5em,
  framexleftmargin=0.5em,
  numbers=none,
}

% ── JuGeo-specific macros ──────────────────────────────────────────────

% Judgment 8-tuple
\newcommand{\J}{\mathcal{J}}
\newcommand{\Jud}[8]{(#1,\,#2,\,#3,\,#4,\,#5,\,#6,\,#7,\,#8)}
\newcommand{\JudShort}[1]{\J_{#1}}

% Components
\newcommand{\coord}{\mathsf{c}}            % coordinate
\newcommand{\prop}{\varphi}                 % proposition
\newcommand{\carrier}{\mathcal{A}}          % carrier type
\newcommand{\evid}{\mathcal{E}}             % evidence bundle
\newcommand{\oblig}{\mathcal{O}}            % obligations
\newcommand{\obstr}{\mathcal{B}}            % obstructions
\newcommand{\trust}{\mathcal{T}}            % trust annotation
\newcommand{\prov}{\Pi}                     % provenance

% Trust algebra
\newcommand{\Talg}{\mathfrak{T}}            % trust ordered algebra
\newcommand{\Eadm}{\mathcal{E}_{\mathrm{adm}}}
\newcommand{\tleq}{\preceq}                 % trust ordering
\newcommand{\tjoin}{\oplus}                 % conservative join
\newcommand{\tatten}{\ominus}               % attenuation
\newcommand{\tpromote}{\uparrow_\pi}        % promotion
\newcommand{\tdemote}{\downarrow_\chi}       % demotion

% Trust tiers
\newcommand{\Tcontra}{\ensuremath{\mathsf{CONTRADICTED}}}
\newcommand{\Tunver}{\ensuremath{\mathsf{UNVERIFIED}}}
\newcommand{\Tcopilot}{\ensuremath{\mathsf{COPILOT}}}
\newcommand{\Toracle}{\ensuremath{\mathsf{ORACLE}}}
\newcommand{\Truntime}{\ensuremath{\mathsf{RUNTIME}}}
\newcommand{\Tsolver}{\ensuremath{\mathsf{SOLVER}}}
\newcommand{\Tproof}{\ensuremath{\mathsf{PROOF}}}

% Site and geometry
\newcommand{\Site}{\mathcal{S}}             % semantic site
\newcommand{\Cov}{\mathrm{Cov}}            % covering family
\newcommand{\Ob}{\mathrm{Ob}}              % objects
\newcommand{\Mor}{\mathrm{Mor}}            % morphisms
\newcommand{\res}{\mathrm{res}}            % restriction
\newcommand{\Cech}{\check{C}}              % Čech complex
\newcommand{\Hcoh}[1]{H^{#1}}             % cohomology group

% Descent
\newcommand{\descent}{\mathrm{desc}}
\newcommand{\glue}{\mathrm{glue}}
\newcommand{\overlap}[2]{#1 \cap #2}

% Semantic moves
\newcommand{\VERIFY}{\textsc{Verify}}
\newcommand{\CONSTRUCT}{\textsc{Construct}}
\newcommand{\NORMALIZE}{\textsc{Normalize}}
\newcommand{\GLUE}{\textsc{Glue}}
\newcommand{\RESTRICT}{\textsc{Restrict}}
\newcommand{\TRANSPORT}{\textsc{Transport}}
\newcommand{\REPAIR}{\textsc{Repair}}
\newcommand{\PROMOTE}{\textsc{Promote}}
\newcommand{\CHALLENGE}{\textsc{Challenge}}
\newcommand{\ESCALATE}{\textsc{Escalate}}

% Fragments
\newcommand{\QFLIA}{\texttt{QF\_LIA}}
\newcommand{\QFLRA}{\texttt{QF\_LRA}}
\newcommand{\QFBV}{\texttt{QF\_BV}}
\newcommand{\QFUF}{\texttt{QF\_UF}}

% Misc
\newcommand{\Python}{\textsc{Python}}
\newcommand{\JuGeo}{\textsc{JuGeo}}
\newcommand{\LEAN}{\textsc{Lean}\,4}
\newcommand{\Fstar}{F$^{\star}$}
\newcommand{\Coq}{\textsc{Coq}}
\newcommand{\Dafny}{\textsc{Dafny}}

% Common shortcuts
\newcommand{\ie}{i.e.\@\xspace}
\newcommand{\eg}{e.g.\@\xspace}
\newcommand{\cf}{cf.\@\xspace}
\newcommand{\wrt}{w.r.t.\@\xspace}

% Evaluation shortcuts
\newcommand{\numcases}{300}
\newcommand{\numspec}{100}
\newcommand{\numeq}{100}
\newcommand{\numbug}{100}
\newcommand{\medtime}{6.5\,ms}
\newcommand{\totaltime}{1.949\,s}


\title{\textbf{Theory-Aware Verification Condition Routing}}

\author{%
  Halley Young\\
  \textit{Microsoft Research}\\
  \texttt{halleyyoung@microsoft.com}
}

\date{}

\begin{document}
\maketitle

\begin{abstract}
Modern verification pipelines send every verification condition (VC)
to a monolithic SMT solver---typically Z3---without first asking
\emph{which decidable fragment the VC belongs to}.
The result is predictable: a quantifier-free linear integer arithmetic
query embedded in a mixed-theory envelope can be 10$\times$ slower than
the same query dispatched to a dedicated decision procedure.
We introduce three components within the \JuGeo{} framework:
\emph{FragmentClassifier}, which maps each VC to the most specific
element of a 14-fragment taxonomy in a mean time of $23.6\,\mu$s;
\emph{FragmentDecomposer}, which applies Nelson--Oppen--style
decomposition to mixed-theory VCs; and
\emph{SolverRouter}, which dispatches each classified fragment to the
backend whose \emph{jurisdiction} covers that fragment, respecting
trust ceilings from the trust algebra of Paper~4.
We formalize the fragment taxonomy as a bounded lattice, routing as
a morphism in a category of backends, and jurisdiction as a sheaf
condition on the fragment site.
An evaluation on 30 VCs across 6 fragment families demonstrates
that classification adds negligible overhead (batch total
$607.5\,\mu$s), that jurisdiction-aware routing directs 26 of 30
VCs to a lightweight arithmetic backend at \$0.001 per VC,
and that \JuGeo{} is the only system among Z3, \Dafny{}/Boogie,
and Why3 that combines fragment classification, multi-backend
routing, jurisdiction checking, trust-aware dispatch, and
Nelson--Oppen decomposition.
\end{abstract}

\tableofcontents
\newpage

%% ═══════════════════════════════════════════════════════════════════════
\section{Introduction}\label{sec:intro}
%% ═══════════════════════════════════════════════════════════════════════

\subsection{The monolithic dispatch problem}

Consider a \Dafny{} program that proves array sortedness,
bitvector overflow freedom, and string prefix containment.
The Boogie intermediate verifier translates each obligation
into an SMT-LIB query and sends \emph{all} of them to a single
invocation of Z3~\cite{z3}.
Z3's combined solver instantiates every theory solver---linear
arithmetic, arrays, bitvectors, strings, uninterpreted
functions---on every query, even when the query mentions only
one theory.
The overhead is measurable: a \QFLIA{} query wrapped in a
\texttt{MIXED} envelope routinely takes 10$\times$ longer
than the same query sent directly to Z3's simplex
engine~\cite{z3}.

The same phenomenon appears in \Fstar{}~\cite{fstar}, which
lowers VCs through a tactic monad into SMT-LIB, and in
Why3~\cite{why3}, which at least allows the user to select
provers via \emph{driver} annotations but does not automate
the selection.
In every case, the tool-chain treats the VC as an opaque
formula and ignores the decidability structure hiding in
plain sight.

\subsection{Our contribution}

This paper introduces \emph{theory-aware verification condition
routing}: the idea that each VC should be classified into the
most specific decidable fragment \emph{before} it is sent to
a solver, and that the router should respect both the
\emph{jurisdiction} of each backend and the \emph{trust ceiling}
imposed by the trust algebra (Paper~4~\cite{paper04}).

Concretely, we contribute:

\begin{enumerate}[label=(\roman*)]
  \item A \textbf{14-fragment taxonomy} organized as a bounded
    lattice, with formal decidability and complexity annotations
    (\Cref{sec:taxonomy}).
  \item \textbf{FragmentClassifier}: a syntactic signature
    extraction algorithm that classifies any VC into the taxonomy
    in $O(|\varphi|)$ time, with soundness and completeness
    guarantees (\Cref{sec:classification}).
  \item \textbf{FragmentDecomposer}: a Nelson--Oppen--style
    decomposition for mixed-theory VCs that produces
    equisatisfiable sub-problems (\Cref{sec:decomposition}).
  \item \textbf{SolverRouter}: a five-strategy dispatcher
    with jurisdiction checking, trust ceilings, and fallback
    chains (\Cref{sec:router}).
  \item A sheaf-theoretic interpretation of jurisdiction as
    a covering condition on the fragment site
    (\Cref{sec:trust-integration}).
\end{enumerate}

\noindent
We evaluate on 30~VCs drawn from six fragment families
(\Cref{sec:evaluation}) and compare against Z3 monolithic,
\Dafny{}/Boogie, and Why3.

%% ═══════════════════════════════════════════════════════════════════════
\section{Background}\label{sec:background}
%% ═══════════════════════════════════════════════════════════════════════

\subsection{SMT-LIB fragments and decidability}

The SMT-LIB standard~\cite{smtlib} organizes first-order theories
into \emph{logics}, each identified by a string such as
\QFLIA{}, \QFLRA{}, or \QFBV{}.
A logic specifies both the background theory and the syntactic
restrictions (quantifier-free, linear, etc.) that guarantee
decidability.
For quantifier-free fragments, decision procedures exist with
known complexity bounds: \QFLIA{} is NP-complete via integer
linear programming~\cite{schrijver}, \QFLRA{} is polynomial
via the simplex method, and \QFBV{} is NP-complete via
bit-blasting~\cite{kroening-strichman}.

\subsection{Nelson--Oppen theory combination}

When a formula mixes symbols from multiple theories
$T_1, \ldots, T_n$, the Nelson--Oppen method~\cite{nelson-oppen}
provides a modular decision procedure, provided each $T_i$
is \emph{stably infinite} and \emph{convex}.
The method partitions the formula into theory-pure sub-formulas,
solves each independently, and propagates equalities over
shared variables until a fixed point is reached.

\begin{definition}[Theory purity]\label{def:purity}
A formula $\varphi$ is \emph{$T_i$-pure} if every
non-logical symbol in~$\varphi$ belongs to the signature
$\Sigma_i$ of~$T_i$, and every variable in~$\varphi$ has a sort
in~$T_i$.
\end{definition}

\subsection{Judgment geometry in brief}

The \JuGeo{} framework (Papers~1--4) organizes verification evidence
into \emph{judgments}
$\J = \Jud{\coord}{\prop}{\carrier}{\evid}{\oblig}{\obstr}{\trust}{\prov}$,
where each component carries semantic meaning: the coordinate~$\coord$
locates the judgment in a semantic site~$\Site$, the
proposition~$\prop$ is the claim, the trust annotation~$\trust$
records provenance, and the evidence bundle~$\evid$ carries
the proof artifact.

The trust algebra $\Talg$ from Paper~4 orders evidence channels:
\[
  \Tcontra \;\tleq\; \Tunver \;\tleq\; \Tcopilot \;\tleq\;
  \Toracle \;\tleq\; \Truntime \;\tleq\; \Tsolver \;\tleq\;
  \Tproof
\]
A backend that returns \texttt{unsat} produces evidence at
level~$\Tsolver$; a backend that times out yields~$\Tunver$;
an LLM oracle yields at most~$\Tcopilot$.
The trust algebra's \emph{no silent promotion} invariant
guarantees that no composition step can silently upgrade
a weak result to a stronger tier.

%% ═══════════════════════════════════════════════════════════════════════
\section{The Fragment Taxonomy}\label{sec:taxonomy}
%% ═══════════════════════════════════════════════════════════════════════

\subsection{The fourteen fragments}

We define a taxonomy of 14 fragments sufficient to cover the VCs
arising in \Dafny{}, \Fstar{}, and \LEAN{} verification
pipelines.

\begin{definition}[Fragment taxonomy]\label{def:taxonomy}
The \emph{fragment taxonomy} is the set
\[
  \mathcal{F} = \bigl\{\,
    \QFLIA,\; \QFLRA,\; \QFBV,\; \QFUF,\;
    \texttt{QF\_AUFLIA},\; \texttt{QF\_ABV},\;
    \texttt{STRINGS},\; \texttt{SEQUENCES},\;
    \texttt{ARRAYS},\; \texttt{DATATYPES},\;
    \texttt{NONLINEAR},\; \texttt{QUANTIFIED},\;
    \texttt{MIXED},\; \texttt{UNKNOWN}
  \,\bigr\}
\]
equipped with a partial order $\sqsubseteq$ defined by
theory inclusion: $F_1 \sqsubseteq F_2$ iff every formula
in~$F_1$ is also in~$F_2$.
\end{definition}

\begin{theorem}[Decidability classification]\label{thm:decidability}
Each fragment $F \in \mathcal{F}$ has a known decidability
status and worst-case complexity:
\begin{center}
\begin{tabular}{@{}lll@{}}
\toprule
\textbf{Fragment} & \textbf{Decidable?} & \textbf{Complexity} \\
\midrule
\QFLIA       & Yes & NP-complete \\
\QFLRA       & Yes & Polynomial (simplex) \\
\QFBV        & Yes & NP-complete (bit-blasting) \\
\QFUF        & Yes & $O(n \log n)$ (congruence closure) \\
\texttt{QF\_AUFLIA} & Yes & NP-complete \\
\texttt{QF\_ABV}    & Yes & NP-complete \\
\texttt{STRINGS}    & Yes & PSPACE \\
\texttt{SEQUENCES}  & Yes & PSPACE \\
\texttt{ARRAYS}     & Yes & NP-complete (QF fragment) \\
\texttt{DATATYPES}  & Yes & NP-complete (QF fragment) \\
\texttt{NONLINEAR}  & Undecidable in general & Semi-decidable \\
\texttt{QUANTIFIED} & Undecidable in general & Semi-decidable \\
\texttt{MIXED}      & Depends on components & Via Nelson--Oppen \\
\texttt{UNKNOWN}    & Unknown & Fallback \\
\bottomrule
\end{tabular}
\end{center}
\end{theorem}

\begin{proof}
The decidability and complexity results for \QFLIA{}, \QFLRA{},
\QFBV{}, and \QFUF{} are classical; see
Kroening and Strichman~\cite{kroening-strichman} for a unified
treatment.
The \texttt{STRINGS} fragment is decidable in PSPACE by the
word-equation results of Makanin~\cite{makanin} as refined by
Plandowski.
\texttt{NONLINEAR} real arithmetic is decidable (Tarski), but
\texttt{NONLINEAR} integer arithmetic is undecidable
(Matiyasevich); we conservatively mark the combined fragment
as undecidable.
\texttt{QUANTIFIED} formulas over arithmetic are undecidable
by the same argument.
\texttt{MIXED} formulas inherit decidability via Nelson--Oppen
when all component theories are decidable, stably infinite,
and convex~\cite{nelson-oppen}.
\end{proof}

\subsection{The fragment lattice}

The partial order $(\mathcal{F}, \sqsubseteq)$ forms a bounded
lattice with bottom element $\QFUF$ (the most restrictive
decidable fragment) and top element $\texttt{UNKNOWN}$.

\begin{figure}[htbp]
\centering
\begin{tikzcd}[row sep=1.8em, column sep=0.6em]
  & & \texttt{UNKNOWN} & & \\
  & \texttt{MIXED} \arrow[ur] & & \texttt{QUANTIFIED} \arrow[ul] & \\
  \texttt{NONLINEAR} \arrow[ur] & \texttt{STRINGS} \arrow[u]
    & \texttt{DATATYPES} \arrow[ul]
    & \texttt{SEQUENCES} \arrow[ull]
    & \texttt{ARRAYS} \arrow[ulll] \\
  & \texttt{QF\_AUFLIA} \arrow[ul] \arrow[u] \arrow[ur]
    & & \texttt{QF\_ABV} \arrow[ul] \arrow[u] & \\
  \QFLIA \arrow[ur] & \QFLRA \arrow[u]
    & \QFBV \arrow[ur] \arrow[uul]
    & & \\
  & & \QFUF \arrow[ull] \arrow[ul] \arrow[u] & &
\end{tikzcd}
\caption{The fragment lattice $(\mathcal{F}, \sqsubseteq)$.
  An arrow $F_1 \to F_2$ means $F_1 \sqsubseteq F_2$, \ie{}
  every formula in~$F_1$ is also in~$F_2$.}
\label{fig:lattice}
\end{figure}

\begin{definition}[Fragment membership]\label{def:membership}
A closed formula $\varphi$ over signature $\Sigma$ belongs to
fragment $F \in \mathcal{F}$, written $\varphi \in F$, if:
\begin{enumerate}[label=(\alph*)]
  \item every non-logical symbol in~$\varphi$ belongs to the
    signature of~$F$,
  \item the syntactic restrictions of~$F$ are satisfied
    (\eg{} quantifier-free, linear, fixed-width bitvector), and
  \item no proper sub-fragment $F' \sqsubset F$ satisfies
    (a) and (b).
\end{enumerate}
Condition~(c) ensures that membership returns the \emph{most
specific} fragment.
\end{definition}

\begin{lemma}[Lattice properties]\label{lem:lattice}
$(\mathcal{F}, \sqsubseteq)$ is a bounded lattice with:
\begin{enumerate}[label=(\roman*)]
  \item Bottom: $\bot = \QFUF$, since uninterpreted functions
    appear in every theory.
  \item Top: $\top = \texttt{UNKNOWN}$, which contains all formulas.
  \item Meet: $F_1 \sqcap F_2$ is the most specific fragment
    containing both $F_1 \cap F_2$.
  \item Join: $F_1 \sqcup F_2$ is the least general fragment
    containing $F_1 \cup F_2$.
\end{enumerate}
\end{lemma}

\begin{proof}
The partial order $\sqsubseteq$ is reflexive, antisymmetric,
and transitive by definition of theory inclusion.
Existence of meets and joins follows from finiteness of
$\mathcal{F}$: for any finite poset, meets and joins exist
if and only if the poset is a lattice, and one verifies
directly that every pair in the Hasse diagram of
\Cref{fig:lattice} has both a least upper bound and a greatest
lower bound.
The bounds $\bot = \QFUF$ and $\top = \texttt{UNKNOWN}$ are
immediate.
\end{proof}

%% ═══════════════════════════════════════════════════════════════════════
\section{Fragment Classification}\label{sec:classification}
%% ═══════════════════════════════════════════════════════════════════════

\subsection{Syntactic signature extraction}

The first step in theory-aware routing is to determine which
fragment a VC belongs to.
We accomplish this via \emph{syntactic signature extraction}:
a single pass over the abstract syntax tree of~$\varphi$
that collects the set of theory symbols, quantifier depth,
and linearity constraints.

\begin{definition}[FragmentClassifier]\label{def:classifier}
The \emph{FragmentClassifier} is a function
\[
  \mathsf{classify} : \mathsf{Formula} \to \mathcal{F}
\]
defined by the following decision tree:
\begin{enumerate}[label=\textbf{Step~\arabic*}.]
  \item Extract the syntactic signature
    $\sigma(\varphi) = (\mathit{sorts},\, \mathit{funs},\,
    \mathit{preds},\, \mathit{qd},\, \mathit{lin})$
    where $\mathit{qd}$ is the quantifier depth and
    $\mathit{lin}$ is a linearity flag.
  \item If $\mathit{qd} > 0$, return $\texttt{QUANTIFIED}$.
  \item If $\sigma$ contains string operations
    (\texttt{str.len}, \texttt{str.++}, \texttt{str.contains}),
    return $\texttt{STRINGS}$.
  \item If $\sigma$ contains bitvector operations
    (\texttt{bvadd}, \texttt{bvand}, \texttt{extract}) and
    array operations (\texttt{select}, \texttt{store}),
    return $\texttt{QF\_ABV}$.
  \item If $\sigma$ contains bitvector operations only,
    return $\QFBV$.
  \item If $\sigma$ contains array operations with integer
    indices, return $\texttt{QF\_AUFLIA}$.
  \item If $\mathit{lin} = \mathsf{false}$ (\ie{} nonlinear
    multiplication), return $\texttt{NONLINEAR}$.
  \item If $\sigma$ contains only integer arithmetic,
    return $\QFLIA$.
  \item If $\sigma$ contains only real arithmetic,
    return $\QFLRA$.
  \item Otherwise, return $\QFUF$.
\end{enumerate}
\end{definition}

\begin{theorem}[Classification soundness]\label{thm:class-sound}
If $\mathsf{classify}(\varphi) = F$, then $\varphi \in F$.
That is, the classifier never assigns a VC to a fragment
whose signature does not contain the VC's symbols.
\end{theorem}

\begin{proof}
We proceed by case analysis on the decision tree.
In each branch, the classifier checks a sufficient condition
for membership: the presence of theory-specific symbols
(\eg{} \texttt{str.len} for \texttt{STRINGS}) or structural
properties ($\mathit{qd} > 0$ for \texttt{QUANTIFIED}).
By \Cref{def:membership}, membership requires that the
formula's symbols belong to the fragment's signature.
Since the decision tree tests exactly this, the assignment
$\mathsf{classify}(\varphi) = F$ implies
$\varphi \in F$.

For the priority ordering: each branch is guarded by the
\emph{absence} of higher-priority symbols.
For instance, \texttt{STRINGS} is checked before \QFLIA{},
so a formula with both string and integer symbols is correctly
classified as \texttt{STRINGS} rather than \QFLIA{}.
\end{proof}

\begin{theorem}[Classification completeness]\label{thm:class-complete}
For every formula $\varphi$ that belongs to some decidable
fragment $F \in \mathcal{F}$, the classifier returns the most
specific fragment: if $\mathsf{classify}(\varphi) = F'$, then
$F' \sqsubseteq F$ for all $F$ with $\varphi \in F$.
\end{theorem}

\begin{proof}
The decision tree is ordered by specificity: quantified
formulas are caught first (the most general undecidable
fragment), then string, bitvector, array, nonlinear, and
finally pure arithmetic and UF.
Within each branch, the test is tight: a formula classified
as \QFLIA{} contains \emph{only} integer linear arithmetic
symbols, so no more specific fragment can contain it.
The base case $\QFUF$ catches every formula not matched by
a more specific branch, and by \Cref{def:membership}(c),
$\QFUF$ is the most specific fragment for such formulas.
\end{proof}

\subsection{Classification in practice}

The following \Python{} listing shows the classifier in action:

\begin{lstlisting}[caption={Fragment classification (simplified).},
  label=lst:classifier]
from enum import Enum
from dataclasses import dataclass

class Fragment(Enum):
    QF_LIA    = "QF_LIA"
    QF_LRA    = "QF_LRA"
    QF_BV     = "QF_BV"
    QF_UF     = "QF_UF"
    QF_AUFLIA = "QF_AUFLIA"
    QF_ABV    = "QF_ABV"
    STRINGS   = "STRINGS"
    SEQUENCES = "SEQUENCES"
    ARRAYS    = "ARRAYS"
    DATATYPES = "DATATYPES"
    NONLINEAR = "NONLINEAR"
    QUANTIFIED = "QUANTIFIED"
    MIXED     = "MIXED"
    UNKNOWN   = "UNKNOWN"

@dataclass(frozen=True)
class SyntacticSignature:
    has_quantifiers: bool
    has_strings: bool
    has_bitvectors: bool
    has_arrays: bool
    has_nonlinear: bool
    has_integers: bool
    has_reals: bool

def classify(sig: SyntacticSignature) -> Fragment:
    if sig.has_quantifiers:
        return Fragment.QUANTIFIED
    if sig.has_strings:
        return Fragment.STRINGS
    if sig.has_bitvectors and sig.has_arrays:
        return Fragment.QF_ABV
    if sig.has_bitvectors:
        return Fragment.QF_BV
    if sig.has_arrays and sig.has_integers:
        return Fragment.QF_AUFLIA
    if sig.has_nonlinear:
        return Fragment.NONLINEAR
    if sig.has_integers:
        return Fragment.QF_LIA
    if sig.has_reals:
        return Fragment.QF_LRA
    return Fragment.QF_UF
\end{lstlisting}

\subsection{Complexity of classification}

\begin{lemma}[Classification cost]\label{lem:class-cost}
$\mathsf{classify}$ runs in $O(|\varphi|)$ time,
where $|\varphi|$ is the size of the formula's AST.
\end{lemma}

\begin{proof}
Signature extraction requires a single depth-first traversal
of the AST, recording the presence of theory-specific symbols
and the quantifier depth.
The decision tree in \Cref{def:classifier} performs a constant
number of comparisons.
Total cost: $O(|\varphi|) + O(1) = O(|\varphi|)$.
\end{proof}


%% ═══════════════════════════════════════════════════════════════════════
\section{Theory Decomposition}\label{sec:decomposition}
%% ═══════════════════════════════════════════════════════════════════════

\subsection{Nelson--Oppen decomposition for MIXED formulas}

When the classifier returns \texttt{MIXED}, the formula spans
multiple theories and cannot be sent to a single specialized
backend.
We apply a Nelson--Oppen--style decomposition~\cite{nelson-oppen}
to split the formula into theory-pure sub-formulas.

\begin{definition}[Theory decomposition]\label{def:decomposition}
Given a formula $\varphi$ over combined signature
$\Sigma = \Sigma_1 \cup \cdots \cup \Sigma_n$,
the \emph{Nelson--Oppen decomposition} produces:
\begin{enumerate}[label=(\alph*)]
  \item Theory-pure sub-formulas
    $\varphi_1, \ldots, \varphi_n$ where each $\varphi_i$
    is $T_i$-pure (\Cref{def:purity}).
  \item A set of shared variables
    $V_{\mathit{shared}} = \bigcup_{i \neq j}
    (\mathsf{vars}(\varphi_i) \cap \mathsf{vars}(\varphi_j))$.
  \item An equality propagation loop that communicates
    entailed equalities $x = y$ for
    $x, y \in V_{\mathit{shared}}$ between the sub-solvers
    until a fixed point.
\end{enumerate}
\end{definition}

\begin{theorem}[Decomposition correctness]\label{thm:decomp}
Let $\varphi$ be a formula over combined signature
$\Sigma_1 \cup \cdots \cup \Sigma_n$ where each $T_i$ is
stably infinite and convex.
Let $\varphi_1, \ldots, \varphi_n$ be the Nelson--Oppen
decomposition of~$\varphi$.
Then:
\[
  \varphi \text{ is satisfiable}
  \;\;\Longleftrightarrow\;\;
  \bigwedge_{i=1}^{n} \varphi_i \text{ is satisfiable
  (under shared-variable propagation)}
\]
\end{theorem}

\begin{proof}
($\Rightarrow$) If $\varphi$ has a model $\mathcal{M}$,
then the restriction $\mathcal{M}|_{\Sigma_i}$ satisfies
$\varphi_i$ for each~$i$, and all shared-variable equalities
holding in $\mathcal{M}$ are consistent across restrictions.

($\Leftarrow$) Suppose each $\varphi_i$ has a model
$\mathcal{M}_i$.
By the stable infiniteness of each~$T_i$, we can extend each
$\mathcal{M}_i$ to have the same cardinality.
By convexity, if $T_i \models \varphi_i \Rightarrow
\bigvee_{k} (x_k = y_k)$, then there exists some~$k$ with
$T_i \models \varphi_i \Rightarrow (x_k = y_k)$, which is
propagated to all other sub-solvers.
The fixed-point iteration terminates because the number of
shared equalities is bounded by $|V_{\mathit{shared}}|^2$,
and each iteration either propagates a new equality or
detects unsatisfiability.
The combined model is obtained by amalgamating the individual
models along the shared variables, which is possible by the
Robinson joint consistency theorem applied to stably infinite
theories~\cite{nelson-oppen}.
\end{proof}

\begin{lemma}[Shared variable propagation]\label{lem:shared-vars}
The equality propagation loop terminates in at most
$\binom{|V_{\mathit{shared}}|}{2}$ rounds, and each
propagated equality preserves satisfiability of the
individual sub-problems.
\end{lemma}

\begin{proof}
Each round either propagates a new equality
$x = y$ with $x, y \in V_{\mathit{shared}}$, refining
the equivalence relation on shared variables, or reaches a
fixed point.
There are at most $\binom{|V_{\mathit{shared}}|}{2}$
distinct equalities, so termination follows.
Preservation of satisfiability holds because a propagated
equality $x = y$ is entailed by $T_i \cup \{\varphi_i\}$;
adding it to $T_j \cup \{\varphi_j\}$ restricts the models
of~$\varphi_j$ but cannot introduce new satisfying assignments,
hence satisfiability is preserved (though the set of models
may shrink).
\end{proof}

\subsection{Decomposition flow}

The following diagram shows the decomposition pipeline:

\begin{figure}[htbp]
\centering
\begin{tikzcd}[row sep=2.5em, column sep=2.8em]
  \varphi_{\texttt{MIXED}}
    \arrow[r, "\text{purify}"]
    \arrow[d, "\text{extract}" swap]
  & \varphi_1 \sqcup \cdots \sqcup \varphi_n
    \arrow[r, "\text{solve}"]
  & \mathsf{Result}_1 \times \cdots \times \mathsf{Result}_n
    \arrow[d, "\text{combine}"] \\
  V_{\mathit{shared}}
    \arrow[r, "\text{propagate}" swap]
  & \text{Equalities}
    \arrow[u, "\text{refine}" near start]
  & \mathsf{Result}_{\texttt{MIXED}}
\end{tikzcd}
\caption{Nelson--Oppen decomposition flow.  The formula
  $\varphi_{\texttt{MIXED}}$ is purified into theory-pure
  sub-formulas, shared variables are extracted, equalities
  are propagated until a fixed point, and individual results
  are combined.}
\label{fig:decomposition}
\end{figure}

\subsection{Decomposition in practice}

\begin{lstlisting}[caption={Theory decomposition (simplified).},
  label=lst:decompose]
@dataclass
class Decomposition:
    sub_formulas: list[tuple[Fragment, str]]
    shared_vars: list[str]

def decompose(formula: str,
              fragment: Fragment) -> Decomposition:
    """Nelson-Oppen decomposition for MIXED formulas."""
    if fragment != Fragment.MIXED:
        return Decomposition(
            sub_formulas=[(fragment, formula)],
            shared_vars=[])
    # Purify: split into theory-pure conjuncts
    conjuncts = extract_conjuncts(formula)
    pure_groups = group_by_theory(conjuncts)
    shared = find_shared_variables(pure_groups)
    # Propagate equalities until fixed point
    propagate_equalities(pure_groups, shared)
    return Decomposition(
        sub_formulas=pure_groups,
        shared_vars=shared)
\end{lstlisting}

%% ═══════════════════════════════════════════════════════════════════════
\section{The SolverRouter}\label{sec:router}
%% ═══════════════════════════════════════════════════════════════════════

\subsection{Backend descriptors}

Each solver backend is described by a \emph{backend descriptor}
that records its capabilities, cost, and trust level.

\begin{definition}[BackendDescriptor]\label{def:backend}
A \emph{backend descriptor} is a tuple
\[
  B = (\mathit{name},\; \mathit{jurisdiction},\;
       \mathit{trust\_ceiling},\; \mathit{cost},\;
       \mathit{latency})
\]
where:
\begin{itemize}
  \item $\mathit{name} \in \mathsf{String}$ identifies the backend.
  \item $\mathit{jurisdiction} \subseteq \mathcal{F}$ is the set
    of fragments the backend can decide.
  \item $\mathit{trust\_ceiling} \in \Eadm$ is the maximum trust
    level the backend can produce.
  \item $\mathit{cost} \in \mathbb{R}_{\geq 0}$ is the per-query
    cost (in monetary units or resource tokens).
  \item $\mathit{latency} \in \mathbb{R}_{\geq 0}$ is the
    expected latency in milliseconds.
\end{itemize}
\end{definition}

\subsection{Routing strategies}

\begin{definition}[Routing strategies]\label{def:strategies}
The \emph{SolverRouter} supports five routing strategies:
\begin{enumerate}[label=(\roman*)]
  \item \textbf{Cheapest}: among backends whose jurisdiction
    covers fragment~$F$, select the one with minimal cost.
  \item \textbf{Fastest}: select the backend with minimal
    expected latency.
  \item \textbf{MostTrusted}: select the backend with the
    highest trust ceiling.
  \item \textbf{RoundRobin}: distribute queries evenly across
    eligible backends.
  \item \textbf{Smart}: a composite strategy that first
    filters by jurisdiction, then selects the cheapest among
    backends within a latency budget.
\end{enumerate}
Each strategy first restricts to backends $B$ with
$F \in B.\mathit{jurisdiction}$; this is the
\emph{jurisdiction guard}.
\end{definition}

\subsection{Routing soundness}

\begin{theorem}[Routing soundness]\label{thm:routing-sound}
For any VC $\varphi$ with $\mathsf{classify}(\varphi) = F$,
the SolverRouter dispatches~$\varphi$ only to a backend~$B$
with $F \in B.\mathit{jurisdiction}$.
\end{theorem}

\begin{proof}
Every routing strategy begins with the jurisdiction guard:
\[
  \mathit{eligible}(F) =
  \{ B \in \mathit{Backends} \mid F \in B.\mathit{jurisdiction} \}
\]
If $\mathit{eligible}(F) = \emptyset$, the router returns a
$\Tunver$ result with a re-dispatch obligation rather than
sending the VC to an incompetent backend.
Otherwise, the strategy selects from $\mathit{eligible}(F)$,
and by construction, the selected backend has $F$ in its
jurisdiction.
\end{proof}

\begin{theorem}[Trust ceiling enforcement]\label{thm:trust-ceiling}
For any VC $\varphi$ dispatched to backend~$B$, the trust
level of the result satisfies
$\trust(\mathsf{result}) \tleq B.\mathit{trust\_ceiling}$.
\end{theorem}

\begin{proof}
The router wraps each backend invocation in a
\emph{trust-capping} combinator:
\[
  \mathsf{cap}(B, r) =
  \begin{cases}
    r & \text{if } \trust(r) \tleq B.\mathit{trust\_ceiling}, \\
    \tatten(r, B.\mathit{trust\_ceiling}) & \text{otherwise}.
  \end{cases}
\]
The attenuation operator $\tatten$ from Paper~4 is monotone
and satisfies $\trust(\tatten(r, c)) \tleq c$ for all~$c$.
Hence $\trust(\mathsf{cap}(B, r)) \tleq B.\mathit{trust\_ceiling}$.
\end{proof}

\subsection{Routing architecture}

\begin{figure}[htbp]
\centering
\begin{tikzcd}[row sep=2.5em, column sep=1.6em]
  & & \mathsf{SolverRouter}
    \arrow[dll, "\text{arith}" swap]
    \arrow[dl, "\text{bv}"]
    \arrow[d, "\text{str}"]
    \arrow[dr, "\text{full}"]
    \arrow[drr, "\text{oracle}"]
    & & \\
  \mathsf{z3\text{-}arith}
  & \mathsf{z3\text{-}bv}
  & \mathsf{z3\text{-}str}
  & \mathsf{z3\text{-}full}
  & \mathsf{copilot}
\end{tikzcd}
\caption{Routing architecture.  The SolverRouter dispatches each
  classified VC to the most appropriate backend.  Arrows represent
  routing channels; each backend has a jurisdiction, trust
  ceiling, and cost.}
\label{fig:routing}
\end{figure}

\subsection{Jurisdiction as a sheaf condition}

We give a categorical interpretation of jurisdiction that
connects to the Grothendieck topology developed in Paper~1.

\begin{definition}[Fragment site]\label{def:fragment-site}
The \emph{fragment site} $\Site_{\mathcal{F}}$ is the category
whose objects are fragments $F \in \mathcal{F}$ and whose
morphisms are inclusions $F_1 \hookrightarrow F_2$ for
$F_1 \sqsubseteq F_2$, equipped with the Grothendieck topology
$\mathcal{J}$ where a covering sieve on~$F$ is any collection
of inclusions $\{F_i \hookrightarrow F\}$ such that
$\bigcup_i F_i = F$.
\end{definition}

A backend descriptor~$B$ defines a presheaf on $\Site_{\mathcal{F}}$:
\[
  \mathcal{B}(F) =
  \begin{cases}
    \{\text{solver instance}\} & \text{if } F \in B.\mathit{jurisdiction}, \\
    \emptyset & \text{otherwise}.
  \end{cases}
\]
The jurisdiction condition states that $B$ can handle a
fragment $F$ if and only if $\mathcal{B}(F) \neq \emptyset$.
The \emph{sheaf condition} for routing requires that compatible
local solutions on an open cover glue to a global solution:

\begin{figure}[htbp]
\centering
\begin{tikzcd}[row sep=2em, column sep=3em]
  \mathcal{B}(F)
    \arrow[r, "\res_{F_1}"]
    \arrow[d, "\res_{F_2}" swap]
  & \mathcal{B}(F_1)
    \arrow[d, "\res_{F_1 \sqcap F_2}"] \\
  \mathcal{B}(F_2)
    \arrow[r, "\res_{F_1 \sqcap F_2}" swap]
  & \mathcal{B}(F_1 \sqcap F_2)
\end{tikzcd}
\caption{Sheaf condition for backend jurisdiction.  If a
  backend covers fragments $F_1$ and $F_2$, its solutions
  must agree on the overlap $F_1 \sqcap F_2$.}
\label{fig:sheaf}
\end{figure}

This ensures that when we decompose a \texttt{MIXED} VC into
sub-fragments and route each to a different backend, the
results are compatible on shared variables---precisely the
condition enforced by Nelson--Oppen propagation.


%% ═══════════════════════════════════════════════════════════════════════
\section{Integration with Trust Algebra}\label{sec:trust-integration}
%% ═══════════════════════════════════════════════════════════════════════

The SolverRouter does not operate in isolation; it is embedded
in the trust algebra $\Talg$ from Paper~4.
Each routing outcome maps to a trust level:

\begin{itemize}
  \item \textbf{Z3 returns \texttt{unsat}}: the VC is
    \emph{solver-discharged}, producing evidence at trust
    level $\Tsolver$.
  \item \textbf{Z3 times out}: the VC is \emph{unverified},
    producing evidence at level $\Tunver$ with a
    re-dispatch obligation attached to the judgment.
  \item \textbf{Copilot oracle responds}: the VC is
    \emph{oracle-suggested}, producing evidence at level
    $\Tcopilot$ with a hard ceiling that prevents
    promotion above $\Tcopilot$ without independent
    verification.
\end{itemize}

\begin{lemma}[Conservative join of routing results]%
\label{lem:conservative-join}
Let $r_1, \ldots, r_n$ be the results of routing the
sub-problems of a decomposed VC $\varphi$ to backends
$B_1, \ldots, B_n$.
The conservative join
$\trust(\bigoplus_{i=1}^{n} r_i) =
\bigsqcap_{i=1}^{n} \trust(r_i)$
preserves the trust ordering: the combined result has
trust level equal to the \emph{minimum} of the individual
trust levels.
\end{lemma}

\begin{proof}
By the definition of the conservative join $\tjoin$
in Paper~4, $\trust(r \tjoin r') = \min(\trust(r), \trust(r'))$.
Induction on~$n$: the base case $n=1$ is trivial.
For $n > 1$, $\trust(\bigoplus_{i=1}^{n} r_i) =
\min(\trust(r_n), \trust(\bigoplus_{i=1}^{n-1} r_i))
= \min(\trust(r_n), \bigsqcap_{i=1}^{n-1} \trust(r_i))
= \bigsqcap_{i=1}^{n} \trust(r_i)$.
This is consistent with the trust algebra's \emph{no silent
promotion} invariant: a chain is only as strong as its
weakest link.
\end{proof}

\begin{lemma}[Re-dispatch obligation]\label{lem:redispatch}
When backend~$B$ times out on VC~$\varphi$ with fragment~$F$,
the router attaches an obligation
$\oblig = \textsc{Re-dispatch}(\varphi, F,
B.\mathit{trust\_ceiling})$ to the judgment.
This obligation can be discharged by:
\begin{enumerate}[label=(\alph*)]
  \item routing to a fallback backend $B'$ with
    $F \in B'.\mathit{jurisdiction}$, or
  \item escalating to the user via the \ESCALATE{} semantic
    move.
\end{enumerate}
\end{lemma}

\begin{proof}
The obligation is a first-class element of the judgment
tuple (\cf{} Paper~1).
Discharge via~(a) follows from \Cref{thm:routing-sound}:
the fallback backend has jurisdiction.
Discharge via~(b) is always available because the \ESCALATE{}
move accepts any obligation and produces an~$\Tunver$ judgment
with human oversight recorded in the provenance~$\prov$.
\end{proof}

The connection to Paper~4's trust soundness theorem is direct:
the SolverRouter is a \emph{trust-monotone} operator.
Formally, if $\trust(\varphi) = \Tunver$ (the VC has not yet
been verified), then after routing and solving:
\[
  \trust(\mathsf{route}(\varphi)) \in
  \{\Tunver,\; \Tcopilot,\; \Toracle,\; \Tsolver\}
\]
and the transition $\Tunver \tleq \trust(\mathsf{route}(\varphi))$
is recorded in the provenance.
No transition can bypass intermediate levels.


%% ═══════════════════════════════════════════════════════════════════════
\section{Evaluation}\label{sec:evaluation}
%% ═══════════════════════════════════════════════════════════════════════

We evaluate the three components---FragmentClassifier,
FragmentDecomposer, and SolverRouter---on a benchmark of 30~VCs
drawn from six fragment families.
All experiments use the \JuGeo{} prototype implementation in
\Python{}.
Total elapsed time for the full pipeline: $0.0023$\,s.

%% ─── §1: Classification ─────────────────────────────────────────────
\subsection{Fragment classification}\label{sec:eval-classify}

We classify all 30~VCs using the FragmentClassifier.
\Cref{tab:fragment-dist} shows the fragment distribution.

\begin{table}[htbp]
\centering
\caption{Fragment distribution across 30 VCs.}
\label{tab:fragment-dist}
\begin{tabular}{@{}lr@{}}
\toprule
\textbf{Fragment} & \textbf{Count} \\
\midrule
\QFUF            & 16 \\
\texttt{STRINGS} & 4  \\
\texttt{QF\_AUFLIA} & 3 \\
\texttt{NONLINEAR}  & 3 \\
\texttt{QUANTIFIED} & 3 \\
\QFBV            & 1  \\
\midrule
\textbf{Total}   & 30 \\
\bottomrule
\end{tabular}
\end{table}

The dominance of \QFUF{} (16 of 30, or 53\%) is consistent
with the observation that many VCs in practice involve only
equality and uninterpreted functions, with arithmetic
appearing in guards and bounds but not in the core obligation.

\Cref{tab:timing} reports classification timing.

\begin{table}[htbp]
\centering
\caption{Classification timing ($\mu$s).}
\label{tab:timing}
\begin{tabular}{@{}lr@{}}
\toprule
\textbf{Metric} & \textbf{Value} \\
\midrule
Mean    & $23.6$  \\
Min     & $16.3$  \\
Max     & $64.7$  \\
Batch total (30 VCs) & $607.5$ \\
\bottomrule
\end{tabular}
\end{table}

The mean classification time of $23.6\,\mu$s per VC is
negligible compared to solver invocation times (typically
milliseconds to seconds).
The maximum of $64.7\,\mu$s corresponds to VC~1
(``array bounds non-negative''), which has the
largest AST in the benchmark.
The batch total of $607.5\,\mu$s for all 30~VCs
confirms that classification is not a bottleneck.

A detailed per-VC breakdown is given in \Cref{tab:per-vc}
(Appendix), but we highlight several representative entries:

\begin{itemize}
  \item VC~11 (``bitwise complement''): the only \QFBV{}
    formula, classified in $21.9\,\mu$s.
  \item VC~6 (``convex combination bound''): classified as
    \texttt{NONLINEAR} in $24.0\,\mu$s due to the presence
    of a nonlinear multiplication $\lambda \cdot x$.
  \item VC~22 (``square non-negative''): classified as
    \texttt{QUANTIFIED} in $21.7\,\mu$s due to the
    universal quantifier $\forall x.\; x^2 \geq 0$.
  \item VC~15 (``array store-select axiom''): classified as
    \texttt{QF\_AUFLIA} in $22.5\,\mu$s due to the
    \texttt{select}/\texttt{store} operations.
\end{itemize}

%% ─── §2: Decomposition ──────────────────────────────────────────────
\subsection{Theory decomposition}\label{sec:eval-decomp}

Four VCs were identified as candidates for decomposition
(VC~26 ``arithmetic + UF'', VC~27 ``bitvector + array'',
VC~28 ``string + arithmetic'', VC~29 ``quantified array'').
In each case, the FragmentClassifier assigned the VC to a
single fragment (\QFUF{}, \texttt{QF\_AUFLIA}, \texttt{STRINGS},
\texttt{QUANTIFIED}, respectively) because the dominant theory
subsumes the secondary theory.
As a result, the decomposer produced single-fragment
decompositions with no shared variables, confirming that
Nelson--Oppen splitting was not needed for these VCs.

This outcome is expected: most VCs in practice are already
theory-pure or near-pure.
The decomposer becomes essential for genuinely mixed VCs, such
as those combining quantified array reasoning with nonlinear
arithmetic---a pattern that arises in loop invariant
verification.

%% ─── §3: Routing ────────────────────────────────────────────────────
\subsection{Solver routing}\label{sec:eval-routing}

With 5~backends registered, the SolverRouter dispatched all
30~VCs according to the Smart strategy.
\Cref{tab:routing} summarizes the routing decisions.

\begin{table}[htbp]
\centering
\caption{Routing decision summary.}
\label{tab:routing}
\begin{tabular}{@{}lrrr@{}}
\toprule
\textbf{Backend} & \textbf{VCs Routed} & \textbf{Cost/VC (\$)} & \textbf{Latency (ms)} \\
\midrule
\texttt{z3-arithmetic} & 26 & 0.001 & 50 \\
\texttt{z3-full}       & 4  & 0.005 & 200 \\
\midrule
\textbf{Total}         & 30 & --- & --- \\
\bottomrule
\end{tabular}
\end{table}

Key observations:
\begin{itemize}
  \item \textbf{Backend utilization}: 26 of 30 VCs (87\%)
    were routed to \texttt{z3-arithmetic}, the lightweight
    backend with the lowest cost (\$0.001/VC) and latency
    (50\,ms).
    Only 4 VCs---all in the \texttt{STRINGS} fragment---required
    the full Z3 combined solver at \$0.005/VC and 200\,ms
    latency.
  \item \textbf{Total cost}: \$0.046 for 30 VCs, with an
    average of \$0.00153 per VC.
  \item \textbf{Latency}: maximum 200\,ms (string VCs),
    average 70\,ms across all VCs.
  \item \textbf{Jurisdiction}: all 30 routing decisions passed
    the jurisdiction check ($\mathit{jurisdiction\_ok} = \mathsf{true}$
    for every VC).
  \item \textbf{Trust ceiling}: all backends reported a trust
    ceiling of $\mathsf{VERIFIED}$, consistent with the
    $\Tsolver$ tier.
  \item \textbf{Fallback chains}: arithmetic-routed VCs had
    fallback chain [\texttt{z3-full}, \texttt{copilot-oracle}];
    string-routed VCs had fallback [\texttt{copilot-oracle}].
\end{itemize}

%% ─── §4: Capability discovery ───────────────────────────────────────
\subsection{Capability discovery}\label{sec:eval-capability}

We tested the backend capability discovery mechanism with
5 queries over different domains.
\Cref{tab:capability} shows the results.

\begin{table}[htbp]
\centering
\caption{Capability discovery results.}
\label{tab:capability}
\begin{tabular}{@{}llr@{}}
\toprule
\textbf{Domain(s)} & \textbf{Capable Backend(s)} & \textbf{Count} \\
\midrule
arithmetic          & z3-arith  & 1 \\
identity            & runtime   & 1 \\
semantic            & copilot   & 1 \\
equality, arithmetic & z3-arith & 1 \\
heap                & (none)    & 0 \\
\bottomrule
\end{tabular}
\end{table}

The ``heap'' domain query returned no capable backends,
demonstrating that the system correctly reports gaps in
coverage rather than routing to an incompetent backend.
This triggers the $\Tunver$ fallback with a re-dispatch
obligation (\Cref{lem:redispatch}).

%% ─── §5: Comparison ────────────────────────────────────────────────
\subsection{Comparison with existing systems}\label{sec:eval-compare}

\Cref{tab:comparison} compares \JuGeo{}'s routing capabilities
with three baselines from the literature.

\begin{table}[htbp]
\centering
\caption{Comparison of VC routing approaches.}
\label{tab:comparison}
\begin{tabular}{@{}lcccc@{}}
\toprule
\textbf{Feature} & \textbf{\JuGeo} & \textbf{Z3}
  & \textbf{\Dafny/Boogie} & \textbf{Why3} \\
\midrule
Fragment classification
  & \checkmark{} (14 fragments)
  & ---
  & Partial (triggers)
  & \checkmark{} (manual) \\
Multi-backend routing
  & \checkmark{} (5 strategies)
  & ---
  & ---
  & \checkmark{} (drivers) \\
Jurisdiction checking
  & \checkmark
  & ---
  & ---
  & --- \\
Trust-aware routing
  & \checkmark
  & ---
  & ---
  & --- \\
Nelson--Oppen decomp.
  & \checkmark
  & Internal
  & ---
  & --- \\
\bottomrule
\end{tabular}
\end{table}

\paragraph{Z3 monolithic~\cite{z3}.}
Z3 uses a combined solver architecture: all VCs are sent to
a single combined theory solver without fragment classification.
Nelson--Oppen combination happens internally but is not
exposed to the user.
There is no notion of jurisdiction, trust, or multi-backend
routing.

\paragraph{\Dafny/Boogie~\cite{dafny,boogie}.}
Boogie generates VCs for Z3, with partial theory awareness
via trigger-based quantifier instantiation.
The user cannot direct specific VCs to specialized backends,
and there is no trust tracking.

\paragraph{Why3~\cite{why3}.}
Why3 supports multiple provers via driver annotations: the user
specifies which prover to use for which theory.
This is the closest to \JuGeo{}'s approach, but the
classification is manual rather than automatic, there is no
jurisdiction checking, and there is no trust algebra integration.

%% ═══════════════════════════════════════════════════════════════════════
\section{Related Work}\label{sec:related}
%% ═══════════════════════════════════════════════════════════════════════

\subsection{SMT-COMP fragment tracks}

The annual SMT-COMP competition~\cite{smtcomp} organizes
benchmarks into fragment tracks (QF\_LIA, QF\_LRA, QF\_BV,
etc.), and competing solvers are evaluated per-track.
This implicitly acknowledges that solvers have varying
performance across fragments.
Our work makes this observation actionable: rather than
submitting every VC to every solver, we classify first and
dispatch to the best backend.

\subsection{Dafny's VC splitting and Boogie}

\Dafny{}~\cite{dafny} and its intermediate verifier
Boogie~\cite{boogie} perform \emph{VC splitting} to manage
the complexity of large verification conditions.
However, this splitting is syntactic (splitting conjunctions)
rather than theory-aware: both halves of a split may contain
the same theories.
Barnett et~al.~\cite{boogie} describe Boogie's architecture
in detail; the VC generator is theory-agnostic.

\subsection{Why3's multi-prover architecture}

Why3~\cite{why3} pioneered multi-prover verification: the user
writes WhyML programs and annotations, and Why3 generates VCs
that can be dispatched to Alt-Ergo, CVC5, Z3, Coq, or
Isabelle via configurable \emph{drivers}.
Filliâtre and Paskevich~\cite{why3} describe the driver
mechanism, which maps WhyML theories to prover-specific
syntax.
This is multi-backend routing, but the routing is manual: the
user must know which prover works best for which theory.
\JuGeo{} automates this decision via fragment classification
and adds jurisdiction and trust guarantees.

\subsection{Z3's combined solver architecture}

Internally, Z3~\cite{z3} uses a variant of Nelson--Oppen
combination with the \emph{model-based theory combination}
(MBTC) approach of de~Moura and Bjørner.
The combined solver instantiates theory solvers lazily, but
the activation decision is made internally and not exposed to
the user.
Our FragmentClassifier can be viewed as an external, pre-solver
version of Z3's internal theory activation logic, with the
added benefit of being able to route to \emph{different}
solver instances.

\subsection{Portfolio solvers}

Portfolio solving~\cite{portfolio} runs multiple solvers in
parallel on the same formula and returns the first result.
This is complementary to our approach: \JuGeo{} classifies
first and dispatches to one backend, but the Smart strategy
could be extended to portfolio mode for high-value VCs.
The key difference is that portfolio solving is
theory-oblivious, while our routing is theory-aware.

\subsection{Satisfiability modulo theories}

The foundational work on SMT solving is due to
Nieuwenhuis, Oliveras, and Tinelli~\cite{smt-survey}, who
formalize the DPLL(T) architecture.
Nelson and Oppen~\cite{nelson-oppen} introduced the theory
combination method that underlies our decomposition.
Barrett et~al.~\cite{smtlib} standardized the SMT-LIB format
and its fragment taxonomy, which we extend with trust and
jurisdiction annotations.

%% ═══════════════════════════════════════════════════════════════════════
\section{Conclusion}\label{sec:conclusion}
%% ═══════════════════════════════════════════════════════════════════════

We have introduced \emph{theory-aware verification condition
routing}, a three-component architecture within the \JuGeo{}
framework that classifies VCs into decidable fragments,
decomposes mixed-theory formulas via Nelson--Oppen, and
dispatches each fragment to the most appropriate backend
with jurisdiction and trust guarantees.

The fragment taxonomy organizes 14 fragments into a bounded
lattice with formal decidability and complexity annotations.
The FragmentClassifier runs in $O(|\varphi|)$ time with a
mean cost of $23.6\,\mu$s per VC, making it negligible
relative to solver invocation.
The SolverRouter supports five strategies and enforces two
critical invariants: routing soundness (jurisdiction coverage)
and trust ceiling enforcement.

Our evaluation on 30~VCs demonstrates that jurisdiction-aware
routing directs 87\% of VCs to a lightweight arithmetic
backend at \$0.001 per VC, with a total cost of \$0.046
for the full benchmark.
A comparison with Z3, \Dafny{}/Boogie, and Why3 shows that
\JuGeo{} is the only system combining automatic fragment
classification, multi-backend routing, jurisdiction checking,
trust-aware dispatch, and Nelson--Oppen decomposition.

The sheaf-theoretic interpretation of jurisdiction as a
covering condition on the fragment site provides a clean
categorical semantics and connects to the Grothendieck
topology of Paper~1.
The trust integration ensures that every routing decision
respects the no-silent-promotion invariant of Paper~4.

Future work includes extending the taxonomy with additional
fragments (sequences, regular expressions, floating-point),
implementing portfolio mode within the Smart strategy, and
integrating with \LEAN{} 4's tactic framework to enable
theory-aware automation within interactive proof assistants.

%% ═══════════════════════════════════════════════════════════════════════
%% BIBLIOGRAPHY
%% ═══════════════════════════════════════════════════════════════════════
\begin{thebibliography}{20}

\bibitem{z3}
L.~de~Moura and N.~Bj\o{}rner.
\newblock Z3: An efficient SMT solver.
\newblock In \emph{TACAS}, LNCS 4963, pp.~337--340. Springer, 2008.

\bibitem{dafny}
K.~R.~M.~Leino.
\newblock Dafny: An automatic program verifier for functional correctness.
\newblock In \emph{LPAR}, LNCS 6355, pp.~348--370. Springer, 2010.

\bibitem{boogie}
M.~Barnett, B.-Y.~E.~Chang, R.~DeLine, B.~Jacobs, and
K.~R.~M.~Leino.
\newblock Boogie: A modular reusable verifier for object-oriented
programs.
\newblock In \emph{FMCO}, LNCS 4111, pp.~364--387. Springer, 2005.

\bibitem{why3}
J.-C.~Filli\^atre and A.~Paskevich.
\newblock Why3 --- Where programs meet provers.
\newblock In \emph{ESOP}, LNCS 7792, pp.~125--128. Springer, 2013.

\bibitem{nelson-oppen}
G.~Nelson and D.~C.~Oppen.
\newblock Simplification by cooperating decision procedures.
\newblock \emph{ACM TOPLAS}, 1(2):245--257, 1979.

\bibitem{smtlib}
C.~Barrett, P.~Fontaine, and C.~Tinelli.
\newblock The SMT-LIB standard: Version 2.6.
\newblock Technical report, 2017.
\newblock \url{https://smtlib.cs.uiowa.edu/}.

\bibitem{fstar}
N.~Swamy, C.~Hri\c{t}cu, C.~Keller, A.~Rastogi, A.~Delignat-Lavaud,
S.~Forest, K.~Bhargavan, C.~Fournet, P.-Y.~Strub,
M.~Kohlweiss, J.-K.~Zinzindohou\'e, and S.~Zanella-B\'eguelin.
\newblock Dependent types and multi-monadic effects in {F*}.
\newblock In \emph{POPL}, pp.~256--270. ACM, 2016.

\bibitem{lean4}
L.~de~Moura and S.~Ullrich.
\newblock The {Lean}~4 theorem prover and programming language.
\newblock In \emph{CADE}, LNAI 12699, pp.~625--635. Springer, 2021.

\bibitem{paper04}
H.~Young.
\newblock Accounting for trust when machines write proofs.
\newblock \emph{Judgment Geometry Series}, Paper~4, 2024.

\bibitem{schrijver}
A.~Schrijver.
\newblock \emph{Theory of Linear and Integer Programming}.
\newblock Wiley, 1986.

\bibitem{kroening-strichman}
D.~Kroening and O.~Strichman.
\newblock \emph{Decision Procedures: An Algorithmic Point of View}.
\newblock Springer, 2nd edition, 2016.

\bibitem{makanin}
G.~S.~Makanin.
\newblock The problem of solvability of equations in a free semigroup.
\newblock \emph{Mat. Sbornik}, 103(2):147--236, 1977.

\bibitem{smtcomp}
SMT-COMP.
\newblock The satisfiability modulo theories competition.
\newblock \url{https://smt-comp.github.io/}, 2023.

\bibitem{smt-survey}
R.~Nieuwenhuis, A.~Oliveras, and C.~Tinelli.
\newblock Solving {SAT} and {SAT} modulo theories: From an abstract
{Davis--Putnam--Logemann--Loveland} procedure to {DPLL(T)}.
\newblock \emph{J. ACM}, 53(6):937--977, 2006.

\bibitem{coq}
The {Coq} Development Team.
\newblock The {Coq} proof assistant.
\newblock \url{https://coq.inria.fr/}, 2023.

\bibitem{portfolio}
Y.~Xu, M.~Stern, and T.~Sandholm.
\newblock Algorithm runtime prediction: Methods and evaluation
(extended abstract).
\newblock In \emph{IJCAI}, pp.~607--613, 2012.

\end{thebibliography}

%% ═══════════════════════════════════════════════════════════════════════
%% APPENDIX
%% ═══════════════════════════════════════════════════════════════════════
\appendix
\section{Per-VC Classification Details}\label{app:per-vc}

\begin{table}[htbp]
\centering
\caption{Complete per-VC classification and routing results.}
\label{tab:per-vc}
\small
\begin{tabular}{@{}rlllrr@{}}
\toprule
\textbf{VC} & \textbf{Description} & \textbf{Fragment}
  & \textbf{Backend} & \textbf{Time ($\mu$s)} & \textbf{Cost (\$)} \\
\midrule
1  & array bounds non-neg.    & \QFUF         & z3-arith & 64.7 & 0.001 \\
2  & loop index in range      & \QFUF         & z3-arith & 30.4 & 0.001 \\
3  & addition commutativity   & \QFUF         & z3-arith & 25.5 & 0.001 \\
4  & scaling preserves sign   & \QFUF         & z3-arith & 26.2 & 0.001 \\
5  & index within length      & \QFUF         & z3-arith & 24.0 & 0.001 \\
6  & convex combination       & \texttt{NL}   & z3-arith & 24.0 & 0.001 \\
7  & average definition       & \QFUF         & z3-arith & 21.1 & 0.001 \\
8  & temperature ceiling      & \QFUF         & z3-arith & 19.0 & 0.001 \\
9  & bvadd commutativity      & \QFUF         & z3-arith & 21.6 & 0.001 \\
10 & unsigned index bound     & \QFUF         & z3-arith & 17.3 & 0.001 \\
11 & bitwise complement       & \QFBV         & z3-arith & 21.9 & 0.001 \\
12 & congruence axiom         & \QFUF         & z3-arith & 21.7 & 0.001 \\
13 & function commutation     & \QFUF         & z3-arith & 19.6 & 0.001 \\
14 & three distinct elems.    & \QFUF         & z3-arith & 17.7 & 0.001 \\
15 & array store-select       & \texttt{AUFL} & z3-arith & 22.5 & 0.001 \\
16 & array read-over-write    & \texttt{AUFL} & z3-arith & 30.4 & 0.001 \\
17 & str length non-neg.      & \texttt{STR}  & z3-full  & 17.7 & 0.005 \\
18 & str concat reflexivity   & \texttt{STR}  & z3-full  & 20.1 & 0.005 \\
19 & substring containment    & \texttt{STR}  & z3-full  & 19.0 & 0.005 \\
20 & constructor test         & \QFUF         & z3-arith & 18.5 & 0.001 \\
21 & head accessor            & \QFUF         & z3-arith & 19.1 & 0.001 \\
22 & square non-neg. ($\forall$)   & \texttt{Q}    & z3-arith & 21.7 & 0.001 \\
23 & perfect square witness   & \texttt{Q}    & z3-arith & 20.5 & 0.001 \\
24 & square non-neg. (NL)     & \texttt{NL}   & z3-arith & 16.3 & 0.001 \\
25 & distance bound           & \texttt{NL}   & z3-arith & 21.2 & 0.001 \\
26 & arithmetic + UF          & \QFUF         & z3-arith & 23.3 & 0.001 \\
27 & bitvector + array        & \texttt{AUFL} & z3-arith & 24.1 & 0.001 \\
28 & string + arithmetic      & \texttt{STR}  & z3-full  & 24.0 & 0.005 \\
29 & quantified array         & \texttt{Q}    & z3-arith & 29.5 & 0.001 \\
30 & pointer identity         & \QFUF         & z3-arith & 24.8 & 0.001 \\
\bottomrule
\end{tabular}

\smallskip
\footnotesize
\texttt{NL} = \texttt{NONLINEAR},\;
\texttt{STR} = \texttt{STRINGS},\;
\texttt{AUFL} = \texttt{QF\_AUFLIA},\;
\texttt{Q} = \texttt{QUANTIFIED}.
\end{table}

\end{document}
